\documentclass[a4paper,12pt]{article}

\usepackage{../../../utils/template}

\newgeometry{ left=0.5cm, right=0.5cm, top=1cm, bottom=1cm }

\begin{document}
    \pagestyle{empty}
    \begin{minipage}[t]{0.30\textwidth}
        \begin{enumerate}[leftmargin=0.2em, itemsep=1.1em]
            \item $\begin{vmatrix}
                       a_{11} & a_{12} & \dots  & a_{1n} \\
                       & a_{22} & \dots  & a_{2n} \\
                       &        & \ddots & \vdots \\
                       &        &        & a_{nn}
            \end{vmatrix}=$
            \item $\begin{vmatrix}
                       a_{11}  & \dots & a_{1,n-1} & a_{1n} \\
                       a_{21} & \dots & a_{2,n-1} & 0      \\
                       \vdots &        & \vdots    & \vdots \\
                       a_{n1} & 0     & \dots & 0
            \end{vmatrix} =$
            \item $\begin{vmatrix}
                       \mathbf{A} & \mathbf{O}          \\
                       \mathbf{O} & \mathbf{B}
            \end{vmatrix} =$
            \item $\begin{vmatrix}
                       \mathbf{O}          & \mathbf{A} \\
                       \mathbf{B} & \mathbf{O}
            \end{vmatrix} =$
            \item $\begin{bmatrix}
                       a_1 &     &     \\
                       & a_2 &     \\
                       &     & a_3
            \end{bmatrix}^n =$
            \item $\begin{bmatrix}
                       O & B \\
                       C & O
            \end{bmatrix}^{-1} =$
        \end{enumerate}
    \end{minipage}
    \hfill
    \begin{minipage}[t]{0.30\textwidth}
        \begin{enumerate}[leftmargin=0.2em, itemsep=1.1em]
            \setcounter{enumi}{6}
            \item $\begin{vmatrix}
                       a_{11} &        &        &        \\
                       a_{21} & a_{22} &        &        \\
                       \vdots & \vdots & \ddots &        \\
                       a_{n1} & a_{n2} & \dots  & a_{nn}
            \end{vmatrix} =$
            \item $\begin{vmatrix}
                       0      & \dots & 0         & a_{1n} \\
                       0      & \dots & a_{2,n-1} & a_{2n} \\
                       \vdots &       & \vdots    & \vdots \\
                       a_{n1} & \dots & a_{n,n-1} & a_{nn}
            \end{vmatrix} =$
            \item $\begin{vmatrix}
                       \mathbf{A} & *          \\
                       \mathbf{O} & \mathbf{B}
            \end{vmatrix} =$
            \item $\begin{vmatrix}
                       *          & \mathbf{A} \\
                       \mathbf{B} & \mathbf{O}
            \end{vmatrix} = $
            \item $\begin{bmatrix}
                       a_1 &     &     \\
                       & a_2 &     \\
                       &     & a_3
            \end{bmatrix}^{-1} =$
            \item $\begin{bmatrix}
                       B & O \\
                       O & C
            \end{bmatrix}^n =$
        \end{enumerate}
    \end{minipage}
    \hfill
    \begin{minipage}[t]{0.30\textwidth}
        \begin{enumerate}[leftmargin=0.2em, itemsep=1.1em]
            \setcounter{enumi}{12}
            \item $\begin{vmatrix}
                       a_{11} &      0  &  \dots      & 0        \\
                       0 & a_{22} &       \dots  &    0    \\
                       \vdots & \vdots & \ddots & \vdots       \\
                       0 & 0 & \dots  & a_{nn}
            \end{vmatrix}=$
            \item $\begin{vmatrix}
                       0      & \dots & 0         & a_{1n} \\
                       0      & \dots & a_{2,n-1} & 0 \\
                       \vdots &       & \vdots    & \vdots \\
                       a_{n1} & \dots & 0 & 0
            \end{vmatrix} =$
            \item $\begin{vmatrix}
                       \mathbf{A} & \mathbf{O} \\
                       *          & \mathbf{B}
            \end{vmatrix} =$
            \item $\begin{vmatrix}
                       \mathbf{O} & \mathbf{A} \\
                       \mathbf{B} & *
            \end{vmatrix} =$
            \item $\begin{bmatrix}
                       &   & a \\
                       & b &   \\
                       c &   &   \\
            \end{bmatrix}^{-1} =$
            \item $\begin{bmatrix}
                       B & O \\
                       O & C
            \end{bmatrix}^{-1} =$
        \end{enumerate}
    \end{minipage}
    \hfill

    \vspace{1em}
    \begingroup
    \renewcommand{\arraystretch}{1.8}
    \setlength{\tabcolsep}{2pt}
    \footnotesize
    \centering
    \resizebox{1\textwidth}{!}{%
        \begin{tabular}{|
                >{\centering\arraybackslash}m{1.6cm}|
                >{\raggedright\arraybackslash}m{3.2cm}|
                >{\raggedright\arraybackslash}m{4cm}|
                >{\raggedright\arraybackslash}m{3.8cm}|
                >{\raggedright\arraybackslash}m{4cm}|
        }
            \hline
            & \multicolumn{1}{c|}{伴随}
            & \multicolumn{1}{c|}{逆}
            & \multicolumn{1}{c|}{转置}
            & \multicolumn{1}{c|}{行列式} \\
            \hline 伴随   & $(A^*)^* =$    & $(A^*)^{-1} =$   & $(A^*)^{T} =$    & $|A^*| =$    \\
            \hline 逆    & $(A^{-1})^* =$ & $(A^{-1})^{-1}=$ & $(A^{-1})^{T} =$ & $|A^{-1}| =$ \\
            \hline 转置   & $(A^T)^* =$    & $(A^T)^{-1} =$   & $(A^T)^T =$      & $|A^T| =$    \\
            \hline 行列式  &                & $|A|^{-1} =$     & $|A|^T =$        & $||A|| =$    \\
            \hline 数    & $(kA)^* =$     & $(kA)^{-1} =$    & $(kA)^T =$       & $|kA| =$     \\
            \hline 穿脱原则 & $(AB)^* =$     & $(AB)^{-1} =$    & $(AB)^T =$       & $|AB| =$     \\
            \hline 加法   &                &                  & $(A + B)^T =$    &              \\
            \hline 交换   & $AA^* =$       & $AA^{-1} =$      &                  &              \\
            \hline
        \end{tabular}%
    }
    \endgroup

    \vspace{1em}
    \begin{minipage}[t]{0.30\textwidth}
        \begin{enumerate}[leftmargin=0.2em, itemsep=1.1em]
            \item $[E_i(k)]^{-1} =$
            \item $E - A^3 =$
            \item $r(A^*) =$
        \end{enumerate}
    \end{minipage}
    \hfill
    \begin{minipage}[t]{0.30\textwidth}
        \begin{enumerate}[leftmargin=0.2em, itemsep=1.1em]
            \setcounter{enumi}{2}
            \item $[E_{ij}(k)]^{-1} =$
            \item $AB - 2B - 4A = 0 \Leftrightarrow$
            \item $|A|$逆序数：
        \end{enumerate}
    \end{minipage}
    \hfill
    \begin{minipage}[t]{0.30\textwidth}
        \begin{enumerate}[leftmargin=0.2em, itemsep=1.1em]
            \setcounter{enumi}{4}
            \item $E_{ij}^{-1} =$
            \item $E + A^3 =$
            \item $\begin{bmatrix}
                       a & b \\
                       c & d
            \end{bmatrix}^{-1} =$
        \end{enumerate}
    \end{minipage}

    \begin{enumerate}
        \setcounter{enumi}{5}
        \item $A^{-1}/A$与$A^*$关系：
        \item 施密特正交化：
    \end{enumerate}

    \newpage

    \begin{minipage}[t]{0.30\textwidth}
        \begin{enumerate}[leftmargin=*, itemsep=1.1em]
            \item $\begin{vmatrix}
                       a_1 & b_1 &        &        & &        \\
                       c_2 & a_2 & b_2    &        & &        \\
                       & c_3 & a_3    & b_3 & &       \\
                       &     & \ddots & \ddots & \ddots &\\
                       &     &        & c_{n-1}& a_{n-1} & b_{n-1} \\
                       &     &        && c_{n}& a_n
            \end{vmatrix} =$
        \end{enumerate}
    \end{minipage}
    \hfill
    \begin{minipage}[t]{0.56\textwidth}
        \begin{enumerate}[leftmargin=*, itemsep=1.1em]
            \setcounter{enumi}{1}
            \item 若$A \sim B$，则
        \end{enumerate}
    \end{minipage}

    \hspace{-0.8em}设$A$是$m \times n$矩阵，$B$满足有关矩阵运算要求的矩阵，则

    \begin{minipage}[t]{1\textwidth}
        \begin{enumerate}[leftmargin=0.2em, itemsep=1.1em]
        \item $r(A)$范围
        \item $r(A)$与哪些矩阵的秩相等
        \item $r(AB)$范围
        \item $r([A, B])$范围
        \end{enumerate}
    \end{minipage}

    \vspace{1em}

    \begin{minipage}[t]{0.48\textwidth}
        \begin{enumerate}[leftmargin=0.2em, itemsep=1.1em]
            \setcounter{enumi}{4}
            \item 若$C$是$s \times t$矩阵，则
            \item 若$A_{m\times n}B_{n\times s} = \mathbf{O}$，则
            \begin{itemize}[leftmargin=*, itemsep=1.0em]
                \item $r(A) + r(B)$
                \item 若 $\mathbf{A}$ 和 $\mathbf{B}$ 为方阵，则
                \item 且$A\;B$非零，则
            \end{itemize}
            \item 对称矩阵
            \item 同解方程组
            \item 线性无关的判定与证明
        \end{enumerate}
    \end{minipage}
    \hfill
    \begin{minipage}[t]{0.52\textwidth}
        \begin{enumerate}[leftmargin=0.2em, itemsep=1.2em]
            \setcounter{enumi}{9}
            \item 齐次方程组有解条件
            \item 正交矩阵
            \item 公共解
            \item 矩阵等价
            \item 矩阵相似
            \item 非齐次方程组有解条件
        \end{enumerate}
    \end{minipage}

    \vspace{2em}

    \renewcommand{\arraystretch}{1.6} % 调整行距
    \begin{tabular} {|
            >{\centering\arraybackslash}m{0.08\textwidth}|
        *{8}{>{\centering\arraybackslash}m{0.08\textwidth}|}
    }
        \hline
        $A$        & $kA+E$ & $A+kE$ & $A^{-1}$ & $A^*$ & $A^T$ & $A^n$ & $P^{-1}AP$ & $f(A)$ \\
        \hline
        $\lambda$  &        &        &          &       &       &       &            &        \\
        \hline
        $\bm{\xi}$ &        &        &          &       &       &       &            &        \\
        \hline
    \end{tabular}

\end{document}
